\newpage
\chapter{問合せ対応}

\begin{enumerate}[label=\thechapter.\arabic*　,leftmargin=*]

  \item \textbf{応対の姿勢}

  問合せへの応対は、学生が活動を始めたり続けたりするための入口である。
  制度の説明にとどまらず、相手が何をしたいのかを汲んで答える。

  教務主事は役割であって身分ではない（学園規則第４章第４条）。
  指導する立場からではなく、対等な立場からの連絡として応対する。

  \item \textbf{窓口}

  問合せは、学園内の連絡手段のほか、学外からは電子メールで受け付ける。

  学外からの問合せには、学生でない者からのものも含まれる。
  入会を検討している者、取材や連携の申入れ、無関係の営業などが想定されるため、相手が誰かを踏まえて応対する。

  \item \textbf{答えられるものと答えられないもの}

  規則やガイドラインに書かれていることは、そのまま答えてよい。
  文書のどこに書かれているかを示すと、次から相手が自分で確認できる。

  一方、まだ決まっていないこと、自分の一存では決められないことは、その場で答えない。
  確認して折り返すと伝えれば足りる。

  答えられない理由が「決まっていないから」であるときは、そのように伝える。
  曖昧な答えでその場をしのぐと、後にそれが前例として扱われる。

  \item \textbf{個別の判断を求められたとき}

  申請が通るかどうか、この行為は違反かといった個別の判断を事前に求められることがある。

  一般的な考え方は答えてよいが、実際の申請や事案を見ずに結論を約束しない。
  審査は申請の内容を見て行うものであり、事前の口頭の回答が審査に優先することはない。

  \item \textbf{相談を受けたとき}

  他の学生の言動についての相談を受けることがある。

  相談は通報とは限らない。
  聞いてほしいだけのこともあるため、本人が何を望んでいるかを確認せずに事案として動かさない。

  ただし、学生の安全に関わる内容であれば、本人の意向にかかわらず対応を検討する。
  その手順は第５章に定める。

  \item \textbf{引き継ぎ}

  自分で判断できない問合せは、他の教務主事に引き継ぐ。

  引き継ぐときは、誰から何を聞かれ、どこまで答えたかを伝える。
  相手に同じ説明を二度させない。

  \item \textbf{記録}

  制度の解釈にわたる問合せと、その回答は記録に残す。

  同じ質問が繰り返されるようであれば、ガイドラインやよくある質問として整備することを検討する。
  問合せ対応は個別の応対であると同時に、文書の不足を知る手がかりでもある。

\end{enumerate}